\begin{landscape}
\thispagestyle{empty}

\section*{ANEXOS}
\addcontentsline{toc}{section}{ANEXOS}

\begin{center}
    \textbf{\Large MATRIZ DE CONSISTENCIA REAJUSTADA} \\
    \vspace{0.2cm}
    \small Título: Efecto del uso de herramientas de inteligencia artificial generativa en la precisión matemática y el rendimiento computacional de simulaciones físicas
\end{center}

\vspace{0.3cm}

\noindent
\begin{tabularx}{\linewidth}{|Y|Y|Y|Y|Y|}
\hline
\rowcolor[gray]{0.9} 
\textbf{Problemas} & \textbf{Objetivos} & \textbf{Hipótesis} & \textbf{Variables} & \textbf{Metodología} \\ \hline

\textbf{Problema General:} \newline
¿Qué efecto tiene el uso de herramientas de IA generativa en la precisión matemática y el rendimiento computacional de simulaciones físicas básicas?  \newline

\textbf{Problemas Específicos:}
\begin{itemize}[leftmargin=*, nosep, itemsep=6pt]
    \item \textbf{PE1:} ¿Qué efecto tiene la generación de código mediante IA en el error cuadrático medio (MSE) de trayectorias físicas? 
    \item \textbf{PE2:} ¿Cómo influye la complejidad del escenario físico generado por IA en el tiempo de ejecución y consumo de memoria? 
    \item \textbf{PE3:} ¿Qué nivel de intervención técnica (correcciones) requiere el código generado para validar su estabilidad numérica? 
\end{itemize} 
& 
\textbf{Objetivo General:} \newline
Evaluar el efecto del uso de herramientas de IA generativa en la precisión matemática y el rendimiento computacional de simulaciones físicas básicas. \newline

\textbf{Objetivos Específicos:}
\begin{itemize}[leftmargin=*, nosep, itemsep=6pt]
    \item \textbf{OE1:} Medir el efecto de la generación de código por IA en el error cuadrático medio de trayectorias físicas respecto al modelo teórico. 
    \item \textbf{OE2:} Determinar el impacto de la complejidad del escenario físico en el rendimiento computacional del código generado. 
    \item \textbf{OE3:} Validar la estabilidad numérica del código generado mediante la cuantificación de las correcciones técnicas requeridas.
\end{itemize} 
& 
\textbf{Hipótesis General:} \newline
El uso de herramientas de IA generativa produce cambios significativos (errores acumulativos e ineficiencias) en la precisión y rendimiento de las simulaciones físicas.  \newline

\textbf{Hipótesis Específicas:}
\begin{itemize}[leftmargin=*, nosep, itemsep=6pt]
    \item \textbf{HIP1:} El código generado por IA presenta un error cuadrático medio significativamente mayor al modelo matemático estándar. 
    \item \textbf{HIP2:} El incremento en la complejidad del escenario físico generado por IA produce un aumento no optimizado en el tiempo de ejecución. 
    \item \textbf{HIP3:} El código de simulación generado por IA requiere intervenciones técnicas manuales para corregir inestabilidades numéricas.
\end{itemize} 
& 
\textbf{Variable Independiente:}
\begin{itemize}[leftmargin=*, nosep]
    \item Herramientas de IA generativa (uso y manipulación de prompts).
\end{itemize}

\textbf{Variables Dependientes:}
\begin{itemize}[leftmargin=*, nosep]
    \item Precisión matemática (Error Cuadrático Medio).
    \item Rendimiento computacional (Tiempo de ejecución y RAM). 
\end{itemize} 
& 
\textbf{Nivel:} Explicativo. \newline
\textbf{Enfoque:} Experimental. \newline
\textbf{Diseño:} Experimental puro (comparativo).  \newline

\textbf{Población/Muestra:} \newline
Algoritmos de simulación física (n=30 ejecuciones por escenario). \newline

\textbf{Técnicas/Instrumentos:} \newline
Experimentación y medición / Hyperfine, GNU time y NumPy.\\ \hline
\end{tabularx}
\end{landscape}